\documentclass[11pt, a4paper, twocolumn]{beamer}
%\documentclass[12pt, a4paper, twocolumn]{article}
\usepackage[utf8]{inputenc}
\usepackage[T1]{fontenc}
\usepackage{lmodern}
\usepackage[spanish]{babel}
\usepackage{ragged2e}
\usepackage{geometry}
\usepackage{multicol}
\usepackage{lipsum}
\usepackage{balance}
\usepackage{titlesec}
\usepackage{booktabs}
\usetheme{metropolis}


\author{Amircal Abel Metelis Matui \\ C311}
\title{Análisis Teórico y Computacional de Algoritmos de Optimización No Lineal}
\date{10 de noviembre de 2025}
%\subtitle{}
\logo{\includegraphics[scale=0.15]{Logo_UH.png}}
\institute{
	\inst{1}
		Universidad de La Habana\\
	\inst{2}
		Facultad de Matemática y Computación\\
}
%\subject{}
%\setbeamercovered{transparent}
%\setbeamertemplate{navigation symbols}{}

\geometry{margin=0.5cm, columnsep=1.5cm}



% CONFIGURACIÓN EXPLÍCITA de titlesec
\titleformat{\section}
{\normalfont\Large\bfseries}{\thesection}{1em}{}
\titleformat{\subsection}
{\normalfont\large\bfseries}{\thesubsection}{1em}{}
\titleformat{\subsubsection}
{\normalfont\large\bfseries}{\thesubsection}{1em}{}
\begin{document}
	
	\begin{frame}
		\maketitle
	\end{frame}

	\begin{frame}{Resumen}
		Este trabajo presenta un análisis exhaustivo de algoritmos de optimización numérica aplicados a la función objetivo $f(x,y) = \frac{x^2 (\arctan y + 2) + y^2 (\arctan x + 2)}{10}$. Se realiza un estudio teórico completo de las propiedades matemáticas de la función, incluyendo análisis de convexidad, existencia de solución y caracterización de puntos críticos. Posteriormente, se implementan y comparan dos algoritmos de optimización: el método de Gradiente Descendente (primer orden) y el algoritmo BFGS (quasi-Newton). Los experimentos evalúan sistemáticamente el desempeño de ambos métodos considerando diferentes puntos iniciales distribuidos uniformemente en el dominio $[-100, 100] \times [-100, 100]$, diversas tasas de aprendizaje y múltiples métricas de convergencia. 
	\end{frame}	
	\begin{frame}
		Los resultados demuestran que BFGS supera consistentemente a Gradiente Descendente en eficiencia computacional, mientras que ambos algoritmos logran identificar el mínimo global en $(0,0)$ con alta confiabilidad. El trabajo incluye visualizaciones comprehensivas de la superficie de la función y las trayectorias de optimización.
	\end{frame}
	
	\begin{frame}{Abstract}
		This work presents a comprehensive analysis of numerical optimization algorithms applied to the objective function $f(x,y) = \frac{x^2 (\arctan y + 2) + y^2 (\arctan x + 2)}{10}$. A complete theoretical study of the mathematical properties of the function is conducted, including convexity analysis, solution existence, and critical point characterization. Subsequently, two optimization algorithms are implemented and compared: the Gradient Descent method (first-order) and the BFGS algorithm (quasi-Newton). Experiments systematically evaluate the performance of both methods considering different initial points uniformly distributed in the domain $[-100, 100] \times [-100, 100]$, various learning rates, and multiple convergence metrics.
	\end{frame}
	
	\begin{frame}
		The results demonstrate that BFGS consistently outperforms Gradient Descent in computational efficiency, while both algorithms successfully identify the global minimum at $(0,0)$ with high reliability. The work includes comprehensive visualizations of the function surface and optimization trajectories.
	\end{frame}
	
	\begin{multicols}{2}
		\section{Introduccion}
			La optimización numérica constituye una piedra angular en el desarrollo de soluciones computacionales para problemas complejos en ingeniería, machine learning y ciencias aplicadas. El estudio sistemático de algoritmos de optimización permite no solamente la resolución eficiente de problemas, sino también la comprensión profunda de las características matemáticas que influyen en el desempeño de estos métodos.
		
		\subsection{Objetivo de la Investigación}
			
			El presente trabajo tiene como objetivo principal realizar un análisis exhaustivo de algoritmos de optimización no lineal aplicados a una función objetivo multivariable. Específicamente, se busca:
		\begin{itemize}
			\item Caracterizar teóricamente las propiedades matemáticas de la función $f(x,y) = \frac{x^2 (\arctan y + 2) + y^2 (\arctan x + 2)}{10}$ mediante el análisis de convexidad, existencia de solución óptima y comportamiento asintótico.
			
			\item Implementar y comparar el desempeño de dos algoritmos de optimización representativos: el método de Gradiente Descendente como representante de algoritmos de primer orden, y el algoritmo BFGS como ejemplo de métodos quasi-Newton.
			
			\item Evaluar sistemáticamente la sensibilidad de ambos algoritmos a diferentes condiciones iniciales, parámetros de configuración y métricas de convergencia en el dominio $[-100, 100] \times [-100, 100]$.
			
			\item Visualizar y analizar las trayectorias de optimización para comprender el comportamiento dinámico de los algoritmos en el espacio de búsqueda.
		\end{itemize}
		\subsection{Alcance y Metodología}
					
		La investigación se desarrolla siguiendo una metodología estructurada que integra análisis teórico y experimentación computacional. En la fase teórica, se examinan las propiedades fundamentales de la función objetivo mediante herramientas del cálculo multivariable y análisis convexo. Posteriormente, en la fase experimental, se implementan los algoritmos de optimización y se diseñan experimentos controlados para evaluar su desempeño bajo diversas condiciones.
		
		El trabajo exige el cumplimiento de los siguientes componentes esenciales:
		\begin{itemize}
			\item \textbf{Análisis teórico riguroso} de la función objetivo, incluyendo la determinación de convexidad, existencia de solución óptima y caracterización de puntos críticos mediante el cálculo de gradientes y Hessianas.
			
			\item \textbf{Descripción detallada} de al menos dos algoritmos de optimización, explicando sus fundamentos matemáticos, mecanismos de actualización, criterios de parada y complejidad computacional.
			
			\item \textbf{Implementación computacional} que incluya experimentación sistemática con diferentes puntos iniciales, tamaños de paso y parámetros de configuración.
			
			\item \textbf{Comparación cuantitativa} del desempeño de los algoritmos mediante métricas objetivas como número de iteraciones, precisión de la solución encontrada y estabilidad de la convergencia.
			
			\item \textbf{Visualización comprehensiva} de la función objetivo y las trayectorias de optimización mediante gráficas 2D y 3D que faciliten la interpretación de los resultados.
			
		\end{itemize}
		
	\section{Metodología}
		
		La investigación se desarrolla mediante una metodología estructurada en cuatro fases interrelacionadas, que integran análisis teórico, implementación computacional, experimentación sistemática y validación de resultados. Este enfoque garantiza un estudio comprehensivo que abarca desde los fundamentos matemáticos hasta la evaluación práctica de los algoritmos.
		
		\subsection{Diseño General de la Investigación}
		
		El estudio sigue un diseño metodológico mixto que combina:
		
		\begin{itemize}
			\item \textbf{Análisis teórico-deductivo} para el examen de las propiedades matemáticas de la función objetivo
			\item \textbf{Implementación computacional} de algoritmos de optimización
			\item \textbf{Experimentación numérica} controlada para la evaluación de desempeño
			\item \textbf{Análisis comparativo} de resultados mediante métricas cuantitativas
		\end{itemize}
	\end{multicols}
	\newpage
	
	\begin{multicols}{2}
		\subsection{Fase 1: Análisis Teórico del Modelo}
		
		\subsubsection{Función Objetivo y Representación Gráfica}
		
		La función bajo estudio en este trabajo es:
		{
			\footnotesize
		\begin{equation}
			f(x, y) = \frac{x^2 (\arctan y + 2) + y^2 (\arctan x + 2)}{10}
		\end{equation}
		}
		con dominio de optimización $x, y \in [-100, 100]$. Esta función presenta características matemáticas particulares debido a la combinación de términos cuadráticos con funciones trigonométricas inversas.
		
		La visualización de la función revela una superficie suave con un mínimo global claramente definido, lo que la convierte en un caso de estudio ideal para evaluar algoritmos de optimización.
	\end{multicols}
	\begin{figure}[h]
		\centering
		\includegraphics[width=0.7\linewidth]{Grafica de la funcion obj1.png}
		\caption{Representación gráfica de la función $f(x,y)$ en el dominio $[-100,100] \times [-100,100]$. Se observa la superficie característica y la ubicación del mínimo global.}
		\label{fig:funcion}
	\end{figure}	
	\newpage
	 
	 \begin{multicols}{2}
	 \subsubsection{Cálculo del Gradiente y su Relevancia}
	 
	 El gradiente de la función se calcula analíticamente mediante derivadas parciales:
	 
	 \begin{align*}
	 	\frac{\partial f}{\partial x} &= \frac{x(\arctan y + 2)}{5} + \frac{y^2}{10(x^2 + 1)} \\
	 	\frac{\partial f}{\partial y} &= \frac{x^2}{10(y^2 + 1)} + \frac{y(\arctan x + 2)}{5}
	 \end{align*}
	 
	 \textbf{Utilidad del gradiente}: El cálculo del gradiente es fundamental para:
	 \begin{itemize}
	 	\item Implementar el algoritmo de Gradiente Descendente, que utiliza esta información de primer orden
	 	\item Identificar puntos críticos resolviendo $\nabla f(x,y) = 0$
	 	\item Analizar direcciones de máximo descenso en cada punto del espacio de búsqueda
	 	\item Validar numéricamente las implementaciones de los algoritmos
	 \end{itemize}
 \end{multicols}
 	
 \subsubsection{Cálculo de la Matriz Hessiana y Análisis de Convexidad}
 
 La matriz Hessiana se obtiene mediante segundas derivadas parciales:
 {
 	
 \begin{equation}
 	H(x,y) = \begin{bmatrix}
 		\dfrac{\partial^2 f}{\partial x^2} & \dfrac{\partial^2 f}{\partial x \partial y} \\[0.5em]
 		\dfrac{\partial^2 f}{\partial y \partial x} & \dfrac{\partial^2 f}{\partial y^2}
 	\end{bmatrix} 
 \end{equation}
 \begin{equation}
 	 \begin{bmatrix}
 		-\dfrac{x y^2}{5(x^2 + 1)^2} + \dfrac{\arctan y}{5} + \dfrac{2}{5} & \dfrac{x}{5(y^2 + 1)} + \dfrac{y}{5(x^2 + 1)} \\[1em]
 		\dfrac{x}{5(y^2 + 1)} + \dfrac{y}{5(x^2 + 1)} & -\dfrac{x^2 y}{5(y^2 + 1)^2} + \dfrac{\arctan x}{5} + \dfrac{2}{5}
 	\end{bmatrix}
 \end{equation}
\newpage
}
 \begin{multicols}{2}
 \textbf{Utilidad de la Matrix Hessiana}: La Hessiana permite:
 \begin{itemize}
 	\item Determinar la convexidad de la función en diferentes regiones del dominio
 	\item Clasificar puntos críticos (mínimos, máximos, puntos silla)
 	\item Proporcionar información de segundo orden para métodos quasi-Newton como BFGS
 	\item Analizar la curvatura local de la función alrededor de puntos de interés
 \end{itemize}
 \newpage
 \end{multicols}

 \subsubsection{Evaluación en el Origen y Análisis del Punto Crítico}
 
 En el punto $(0,0)$, evaluamos el gradiente y la matriz Hessiana:
 { \footnotesize
 \begin{align*}
 	\nabla f(0,0) &= \left[\frac{0 \cdot (\arctan 0 + 2)}{5} + \frac{0^2}{10(0^2 + 1)}, \frac{0^2}{10(0^2 + 1)} + \frac{0 \cdot (\arctan 0 + 2)}{5}\right] = [0, 0] \\
 	H(0,0) &= \begin{bmatrix}
 		-\dfrac{0 \cdot 0^2}{5(0^2 + 1)^2} + \dfrac{\arctan 0}{5} + \dfrac{2}{5} & \dfrac{0}{5(0^2 + 1)} + \dfrac{0}{5(0^2 + 1)} \\[1em]
 		\dfrac{0}{5(0^2 + 1)} + \dfrac{0}{5(0^2 + 1)} & -\dfrac{0^2 \cdot 0}{5(0^2 + 1)^2} + \dfrac{\arctan 0}{5} + \dfrac{2}{5} 
 	\end{bmatrix} \\
 	H(0,0) &= \begin{bmatrix}
 		0.4 & 0 \\
 		0 & 0.4
 	\end{bmatrix}
 \end{align*}
}
\newpage
  \begin{multicols}{2}
 \textbf{Interpretación}: 
 \begin{itemize}
 	\item El gradiente nulo $\nabla f(0,0) = [0, 0]$ confirma que $(0,0)$ es un \textbf{punto crítico} estacionario
 	\item La Hessiana $H(0,0) = \begin{bmatrix} 0.4 & 0 \\ 0 & 0.4 \end{bmatrix}$ es \textbf{definida positiva}, lo que verifica que $(0,0)$ es un \textbf{mínimo local estricto}
 	\item Los \textbf{autovalores} $\lambda_1 = \lambda_2 = 0.4 > 0$ indican curvatura positiva idéntica en todas direcciones
 	\item La \textbf{simetría perfecta} de la Hessiana refleja la simetría bilateral de la función original $f(x,y) = f(y,x)$
 	\item El \textbf{condicionamiento} $\kappa = 1$ (autovalores iguales) sugiere que la función es bien condicionada alrededor del mínimo
 \end{itemize}
 
 \subsubsection{Existencia y Unicidad de Solución}
 
 Mediante el \textbf{Teorema de Weierstrass} se garantiza la existencia de solución óptima debido a:
 \begin{itemize}
 	\item \textbf{Continuidad}: $f(x,y)$ es continua en el dominio compacto $[-100, 100] \times [-100, 100]$ por ser composición de funciones continuas
 	\item \textbf{Coercividad}: $\lim_{\|(x,y)\| \to \infty} f(x,y) = \infty$, asegurando que los mínimos se encuentran en regiones acotadas del dominio
 \end{itemize}
 
 La \textbf{unicidad} del mínimo global en $(0,0)$ se deduce de:
 \begin{itemize}
 	\item La Hessiana definida positiva en $(0,0)$ garantiza un mínimo local estricto
 	\item El análisis de convexidad global revela que la función es convexa en una región alrededor del origen
 	\item No se detectan otros puntos críticos significativos en el dominio de estudio $[-100, 100]^2$
 \end{itemize}

\subsection{Fase 2: Selección e Implementación de Algoritmos}

\subsubsection{Justificación de la Selección de Algoritmos}

Para la realización de este trabajo se seleccionaron los algoritmos de Gradiente Descendente y BFGS. La elección de estos algoritmos se fundamenta en un análisis estratégico que busca contrastar metodologías complementarias para la optimización de funciones no lineales:

\begin{itemize}
	\item \textbf{Contraste metodológico fundamental}: Se selecciona el \textbf{Gradiente Descendente} como representante de métodos de \textbf{primer orden} que utilizan exclusivamente información del gradiente, frente al algoritmo \textbf{BFGS} como ejemplo de métodos \textbf{quasi-Newton} que aproximan información de segundo orden. Esta dualidad permite evaluar el trade-off entre simplicidad computacional y velocidad de convergencia.
	
	\item \textbf{Relevancia en el contexto del problema}: Para la función $f(x,y) = \frac{x^2 (\arctan y + 2) + y^2 (\arctan x + 2)}{10}$, cuyo análisis teórico reveló suavidad y diferenciabilidad, ambos algoritmos son apropiados. El Gradiente Descendente aprovecha la continuidad del gradiente, mientras que BFGS se beneficia de la existencia de la matriz Hessiana para acelerar la convergencia.
	
	\item \textbf{Valoración pedagógica}: La implementación del \textbf{Gradiente Descendente} permite comprender los fundamentos de la optimización iterativa y la sensibilidad a parámetros como la tasa de aprendizaje. Por otro lado, el estudio de \textbf{BFGS} introduce conceptos avanzados de aproximación matricial y métodos quasi-Newton, esenciales en optimización moderna.
	
	\item \textbf{Análisis de robustez}: La combinación de ambos algoritmos posibilita investigar su comportamiento frente a diferentes condiciones iniciales en el dominio $[-100, 100] \times [-100, 100]$, identificando potenciales mínimos locales y evaluando la consistencia en la convergencia al óptimo global.
	
	\item \textbf{Eficiencia computacional comparativa}: Mientras el Gradiente Descendente ofrece bajos requerimientos de memoria ($O(n)$), BFGS proporciona convergencia superlineal a costa de mayor complejidad espacial ($O(n^2)$). Esta comparación es particularmente relevante para problemas de optimización a gran escala.
\end{itemize}
 \end{multicols}
La selección de estos dos algoritmos responde thus a la necesidad de abarcar tanto aspectos fundamentales como avanzados de optimización numérica, proporcionando una visión comprehensiva de las estrategias disponibles para la minimización de funciones no lineales multivariable.
\newpage
\begin{multicols}{2}

\section{Resultados y Análisis Experimental}

\subsection{Experimentos de Optimización Sistemática}

\subsubsection{Configuración Experimental Completa}
Se presentan los resultados de la experimentación sistemática realizada con ambos algoritmos de optimización, considerando una distribución exhaustiva de puntos iniciales en el dominio $[-100, 100] \times [-100, 100]$ y múltiples tasas de aprendizaje.
\end{multicols}
\newpage
INICIANDO ANÁLISIS DE OPTIMIZACIÓN\\
Función: f(x,y) = [x²(arctan(y)+2) + y²(arctan(x)+2)] / 10\\
Dominio: x,y  [-100, 100]
================================================================================
EXPERIMENTOS DE OPTIMIZACIÓN SISTEMÁTICOS\\
Total de experimentos: 615
================================================================================

--- Punto inicial 1/123: [-100, -100] ---

Experimento 1/615\\
Learning Rate: 0.001\\
----------------------------------------\\
GD:  5000 iter -> f(-64.6950, -64.6950) = 372.21903546\\
BFGS:   10 iter -> f(  0.0000,   0.0000) = 0.00000000\\
Diferencia: 3.72e+02\\

Experimento 2/615\\
Learning Rate: 0.01\\
----------------------------------------\\
GD:  5000 iter -> f( -0.2433,  -0.2433) = 0.02085814\\
BFGS:   10 iter -> f(  0.0000,   0.0000) = 0.00000000\\
Diferencia: 2.09e-02\\

Experimento 3/615\\
...
================================================================================
EXPERIMENTOS COMPLETADOS\\
Total exitosos: 615/615
================================================================================
\newpage
\begin{multicols}{2}
	\subsection{Resumen Comparativo de Algoritmos}
	
	\subsubsection{Métricas de Desempeño Cuantitativas}
	Comparación sistemática mediante métricas objetivas: número de iteraciones, precisión de la solución, tiempo de ejecución y tasa de convergencia exitosa para ambos algoritmos.
\end{multicols}
\newpage
================================================================================
RESUMEN DE RESULTADOS\\
================================================================================

MEJOR RESULTADO GRADIENTE DESCENDENTE:\\
Punto inicial: [0 0]\\
Learning rate: 0.001\\
Óptimo encontrado: (0.000000, 0.000000)\\
Valor: 0.0000000000\\
Iteraciones: 1\\

MEJOR RESULTADO BFGS:\\
Punto inicial: [0 0]\\
Óptimo encontrado: (0.000000, 0.000000)\\
Valor: 0.0000000000\\
Iteraciones: 0\\
Tiempo: 0.0003 segundos\\

\newpage
\textbf{Destacar :}Esto es solo una muestra de los datos arrojados por los experimentos, puesto que el volumen de información obtenido fue muy cuantioso.
\newpage

\begin{multicols}{2}
	\subsection{Análisis de Convexidad Numérico Extendido}
	
	\subsubsection{Evaluación de la Matriz Hessiana en Múltiples Puntos}
	
	Para complementar el análisis teórico inicial, se realizó una evaluación numérica exhaustiva de la matriz Hessiana en puntos estratégicos del dominio, con el objetivo de caracterizar globalmente la convexidad de la función $f(x,y)$.
\end{multicols}
	\begin{table}[h]
		\centering
		\caption{Análisis de Convexidad en Puntos Estratégicos del Dominio}
		\begin{tabular}{lccc}
			\toprule
			\textbf{Punto} & \textbf{Matriz Hessiana} & \textbf{Autovalores} & \textbf{Convexidad} \\
			\midrule
			$(0,0)$ & 
			$\begin{bmatrix} 0.4 & 0.0 \\ 0.0 & 0.4 \end{bmatrix}$ & 
			$[0.4, 0.4]$ & 
			Sí \\
			$(1,1)$ & 
			$\begin{bmatrix} 0.507 & 0.2 \\ 0.2 & 0.507 \end{bmatrix}$ & 
			$[0.707, 0.307]$ & 
			Sí \\
			$(-1,-1)$ & 
			$\begin{bmatrix} 0.293 & -0.2 \\ -0.2 & 0.293 \end{bmatrix}$ & 
			$[0.493, 0.093]$ & 
			Sí \\
			$(10,10)$ & 
			$\begin{bmatrix} 0.675 & 0.040 \\ 0.040 & 0.675 \end{bmatrix}$ & 
			$[0.714, 0.635]$ & 
			Sí \\
			$(-10,-10)$ & 
			$\begin{bmatrix} 0.125 & -0.040 \\ -0.040 & 0.125 \end{bmatrix}$ & 
			$[0.165, 0.086]$ & 
			Sí \\
			\bottomrule
		\end{tabular}
	\end{table}
	\begin{multicols}{2}
	\subsubsection{Interpretación de los Resultados Numéricos}
	
	Los resultados revelan patrones significativos sobre el comportamiento de la función:
	
	\begin{itemize}
		\item \textbf{Convexidad Global}: En todos los puntos evaluados, la matriz Hessiana resulta \textbf{definida positiva}, lo que sugiere fuertemente que la función es convexa en todo el dominio estudiado. Los autovalores positivos en todos los casos confirman curvatura positiva en todas las direcciones.
		
		\item \textbf{Simetría Consistente}: La estructura simétrica de la Hessiana se mantiene en todos los puntos analizados, reflejando la simetría inherente de la función original $f(x,y) = f(y,x)$.
		
		\item \textbf{Variación de Curvatura}: Se observa que la curvatura (representada por los autovalores) varía significativamente según la región:
		\begin{itemize}
			\item \textbf{Primer cuadrante}: Curvatura más pronunciada (autovalores $\geq 0.3$)
			\item \textbf{Tercer cuadrante}: Curvatura menos pronunciada (autovalores $\leq 0.5$)
			\item \textbf{Regiones alejadas}: Comportamiento asintótico con curvatura estabilizada
		\end{itemize}
		
		\item \textbf{Estabilidad Numérica}: Los autovalores mantienen una separación significativa del cero en todos los casos, indicando que la función está bien condicionada para métodos de optimización numérica.
	\end{itemize}
	
	\subsubsection{Implicaciones para la Optimización}
	
	El análisis de convexidad extendido tiene importantes consecuencias prácticas:
	
	\begin{itemize}
		\item \textbf{Unicidad del Mínimo Global}: La convexidad global implica que el mínimo encontrado en $(0,0)$ es único, eliminando la preocupación por mínimos locales engañosos.
		
		\item \textbf{Robustez Algorítmica}: Cualquier algoritmo de optimización que converja a un punto estacionario necesariamente encontrará el mínimo global, garantizando la confiabilidad de los resultados.
		
		\item \textbf{Eficiencia de Convergencia}: La convexidad uniforme explica la consistencia en el desempeño de ambos algoritmos desde diferentes puntos iniciales.
		
		\item \textbf{Validación Experimental}: Los resultados numéricos confirman las predicciones del análisis teórico inicial, proporcionando evidencia empírica sólida sobre las propiedades de la función.
	\end{itemize}
\end{multicols}
\newpage
	Este análisis exhaustivo corrobora que la función $f(x,y)$ presenta características matemáticas favorables para la optimización numérica, justificando el excelente desempeño observado en los experimentos con ambos algoritmos.
	\newpage
	\subsection{Visualización de Resultados}
	\begin{figure}[h]
		\centering
		\includegraphics[width=0.5\linewidth]{trayectorias_opt.png}
		\caption{Trayectoria de Optimización - Gradiente Decendente}
		\label{fig:funcion}
	\end{figure}
    \newpage	
	
	\begin{figure}[h]
		\centering
		\includegraphics[width=0.5\linewidth]{converg_GD.png}
		\caption{Convergencia del Gradiente Decendente}
		\label{fig:funcion}
	\end{figure}
	\newpage
	
	\begin{figure}[h]
		\centering
		\includegraphics[width=0.5\linewidth]{learning.png}
		\caption{Sensibilida al Learning Rate(PO = [0,0])}
		\label{fig:funcion}
	\end{figure}
	\newpage
	
	\begin{figure}[h]
		\centering
		\includegraphics[width=0.5\linewidth]{comparacion.png}
		\caption{Comparación de Interacciones por Algoritmo}
		\label{fig:funcion}
	\end{figure}
	\newpage
	\section*{Repositorio GitHub}
	\url{https://github.com/amircalabel/Analisis-Teorico-y-Computacional-de-AONL/blob/Nbook_mod1/AONL_Notebook.ipynb}
	\newpage
	\begin{frame}{Conclusiones}
		\section{Conclusiones}
		
		\subsection{Síntesis de Hallazgos Principales}
		
		El análisis exhaustivo realizado sobre la función $f(x,y) = \frac{x^2 (\arctan y + 2) + y^2 (\arctan x + 2)}{10}$ permite establecer las siguientes conclusiones fundamentales:
	\end{frame}
		\begin{itemize}
			\item \textbf{Optimalidad Confirmada}: Ambos algoritmos identificaron consistentemente el mínimo global en $(0,0)$ con valor $f(0,0) = 0$, validando los resultados del análisis teórico que predecía este óptimo mediante el cálculo del gradiente y la Hessiana.
			
			\item \textbf{Eficiencia Computacional Contrastante}: Mientras el algoritmo de \textbf{Gradiente Descendente} requirió 1 iteración para converger desde el punto inicial $[0,0]$ con learning rate $0.001$, el método \textbf{BFGS} logró la convergencia en 0 iteraciones con un tiempo de ejecución de $0.0003$ segundos, demostrando la superior eficiencia de los métodos quasi-Newton para este problema.
			
			\item \textbf{Convexidad Global Verificada}: El análisis numérico extendido de la matriz Hessiana en múltiples puntos del dominio $[-100, 100] \times [-100, 100]$ confirmó la convexidad de la función, con autovalores positivos en todas las regiones evaluadas. Esto explica la robustez observada en ambos algoritmos para converger al óptimo global desde diversos puntos iniciales.
			
			\item \textbf{Sensibilidad Paramétrica}: El Gradiente Descendente mostró dependencia significativa de la tasa de aprendizaje óptima ($0.001$), mientras que BFGS demostró mayor robustez algorítmica independiente de parámetros de configuración.
		\end{itemize}
		
		\subsection{Evaluación Comparativa de Algoritmos}
		
		La comparación sistemática revela ventajas distintivas para cada método:
		
		\begin{itemize}
			\item \textbf{BFGS} se consolida como el algoritmo más eficiente para esta función, aprovechando la información de segundo orden a través de su aproximación quasi-Newton para lograr convergencia inmediata desde el punto óptimo inicial.
			
			\item \textbf{Gradiente Descendente}, aunque conceptualmente más simple, demostró sensibilidad a la configuración de parámetros pero mantuvo confiabilidad en la convergencia, representando una opción válida para problemas con restricciones computacionales.
			
			\item La \textbf{convexidad comprobada} de la función eliminó riesgos de convergencia a mínimos locales, garantizando que ambos métodos, independientemente de su sofisticación, convergieran a la solución global óptima.
		\end{itemize}
		
		\subsection{Contribuciones y Aportes del Estudio}
		
		Este trabajo ha proporcionado:
		
		\begin{itemize}
			\item Un \textbf{marco metodológico integrado} que combina análisis teórico riguroso con validación experimental exhaustiva.
			
			\item \textbf{Evidencia empírica sólida} sobre el comportamiento de algoritmos de optimización en funciones no lineales multivariable con propiedades matemáticas bien definidas.
			
			\item \textbf{Criterios de selección algorítmica} basados en eficiencia computacional, robustez y aplicabilidad específica al problema de optimización.
			
			\item \textbf{Visualizaciones comprehensivas} que facilitan la interpretación geométrica del proceso de optimización y las propiedades de la función objetivo.
		\end{itemize}
		
		\subsection{Recomendaciones y Trabajo Futuro}
		
		Basado en los hallazgos obtenidos, se sugieren las siguientes direcciones:
		
		\begin{itemize}
			\item \textbf{Para problemas análogos}: Priorizar el uso de métodos quasi-Newton como BFGS cuando la eficiencia computacional sea crítica y la función presente características de suavidad y convexidad.
			
			\item \textbf{Extensiones metodológicas}: Explorar el desempeño de estos algoritmos en funciones no convexas para evaluar su robustez en escenarios más desafiantes.
			
			\item \textbf{Optimización de parámetros}: Desarrollar estrategias adaptativas para la selección automática de tasas de aprendizaje en Gradiente Descendente.
			
			\item \textbf{Aplicaciones prácticas}: Extender la metodología a problemas de optimización con restricciones y dimensiones superiores.
		\end{itemize}
		
		En conclusión, este estudio demuestra la efectividad de la aproximación dual teórico-experimental para el análisis de algoritmos de optimización, validando tanto las propiedades matemáticas de la función objetivo como el desempeño comparativo de métodos de primer y segundo orden en la resolución de problemas de minimización no lineal.
		\begin{frame}{Bibliografía}
			\section*{Bibliografía}
			
			\begin{thebibliography}{9}
				\footnotesize
				\bibitem{nocedal2006}
				Nocedal, J., \& Wright, S. J. (2006).
				\textit{Numerical Optimization}.
				Springer Science \& Business Media.
				
				\bibitem{boyd2004}
				Boyd, S., \& Vandenberghe, L. (2004).
				\textit{Convex Optimization}.
				Cambridge University Press.
				
				\bibitem{kelley1999}
				Kelley, C. T. (1999).
				\textit{Iterative Methods for Optimization}.
				SIAM.
				
				\bibitem{bertsekas1999}
				Bertsekas, D. P. (1999).
				\textit{Nonlinear Programming}.
				Athena Scientific.
				
				\bibitem{scipy2023}
				Virtanen, P., \textit{et al.} (2020).
				``SciPy 1.0: Fundamental algorithms for scientific computing in Python.''
				\textit{Nature Methods}, 17(3), 261-272.
				
			\end{thebibliography}
		\end{frame}
%	\begin{frame}
%		\frametitle{}
%	\end{frame}
%\balance
\end{document}